\begin{table}[t]
\centering
\caption{Stepwise simulation workflow used for the baseline exact-position rate calculation. The table is a manuscript-facing compression of the validated workflow; detailed validation files retain the run-level details.}
\label{tab:simulation_workflow}
\begingroup
\setlength{\tabcolsep}{2pt}
\renewcommand{\arraystretch}{0.78}
\tiny
\begin{tabular}{@{}p{0.12\linewidth}p{0.25\linewidth}p{0.36\linewidth}p{0.18\linewidth}@{}}
\toprule
Step & Role in the calculation & Current implementation or evidence & Boundary for interpretation \\
\midrule
00--01 Geometry & Define detector mass model, active volumes, side-entry frame, and injection surfaces & Baseline detector/cryostat geometry: TES array, Be side window, Cu/W passive material, segmented CsI shield; window tilted $45^\circ$ upward & Geometry closure only; not a rate calculation \\
02 Prompt and buildup transport & Generate prompt atmospheric events and activation histories & Detector/cryostat prompt and buildup streams, each with 25,210,216 particles; prompt rates use per-species exposure normalization & Low-stat branches are closure checks; upstream optics mass is outside the primary rate budget \\
03 Delayed activity and source & Convert activation products to a fixed day-15 delayed source & Detector/cryostat delayed activity $85.636696\,\mathrm{Bq}$; nominal exact-position point-source support with $M=50000$ and $10^6$ delayed triggers & Activity preserving; Step10 tests convergence; upstream Ge-proxy delayed activity is a separate side component \\
04 Laue optics & Compute 511 keV effective area and focused focal-plane phase space & Ge(111), 10 m optical Monte Carlo from XOP/CRYSTAL rocking curves; $\aeff(511)=20.08476\,\mathrm{cm^2}$ and spot contained by Be aperture & Signal only; focal crossings are transported 511 keV gammas, not optics self-background \\
09 Optics-detector coupling & Replay tracked optical focal crossings in the detector & Accepted focal gamma rays are written as a MEGAlib particle-list source and transported through the baseline detector setup & Tracked focal crossings only; signal replay does not represent optics self-background \\
Aux. optics hardware & Check upstream-lens prompt/activation boundary & Equal-volume active-Ge lens proxy at 10 m; far-field source radius $1060\,\mathrm{cm}$ from a $1008.8\,\mathrm{cm}$ detector-plus-lens bound; eight EXPACS/PARMA families close prompt and activation-production transport; isolated Ge-proxy delayed transport selects zero W2 events, with a 95\% upper rate $1.25\times10^{-4}\cps$ & Delayed Ge-proxy side component closed; prompt self-background isolation and support hardware remain outside the primary budget \\
05 Detector response & Apply common selection to prompt, delayed, and focused streams & Active-shield post-processing veto, side-entry Compton/FoV consistency, and $\wii$ hard window; day-15 $\wii$ rates $B/S=0.0629804/0.00118117\cps$ & Veto defined by post-processing selection, not production-time trigger filtering \\
06 Mission-time fold & Propagate selected rates through balloon mission time & Time-dependent prompt, delayed activity, atmosphere, and accidental-veto live-factor fold; mean $\wii$ rates $B/S=0.0631923/0.00117281\cps$ & Rate-level fold; no new transport at each time bin \\
07 Source cases & Map detector response to source-flux units & W2 response $11.8117\,\mathrm{cps}/(\mathrm{ph\,cm^{-2}\,s^{-1}})$ for the reference point-source case & Point-source normalization kept separate from diffuse-sky checks \\
08 Significance & Compute cumulative significance and flux thresholds & Hard-window result: $Z_{20\mathrm{d}}=6.13039$, $T_{3\sigma}=4.77766\,\mathrm{d}$, $\fthree(20\mathrm{d})=4.89365\times10^{-5}\cms$ & Gaussian counting metric for consistent comparisons \\
10 Robustness checks & Test support-size and seed dependence at selected-rate level & Four transport-backed cases over $M=5000$, 20,000, and 50,000; W2 background and $Z_{20\mathrm{d}}$ ranges are 1.119\% and 0.551\% & Supports exact-position delayed transport for the primary rates \\
\bottomrule
\end{tabular}
\endgroup
\end{table}
